% Package dan pengaturan umum. Konfigurasi/cover ada di main.tex.
\usepackage[utf8]{inputenc}
\usepackage[T1]{fontenc}
\usepackage[indonesian]{babel}
\usepackage[margin=2.5cm]{geometry}
\usepackage{graphicx}
\usepackage{booktabs}
\usepackage{longtable}
\usepackage{array}
\usepackage{amsmath,amssymb}
\usepackage{listings}
\usepackage{xcolor}
\usepackage{float}
\usepackage{caption}
\usepackage{subcaption}
\usepackage{tikz}
\usepackage{setspace}
\usepackage{hyperref}
\usepackage{fancyhdr}

\graphicspath{{figures/}}
\setstretch{1.15}
\hypersetup{colorlinks=true, linkcolor=black, urlcolor=blue, citecolor=black}

\lstset{
  basicstyle=\ttfamily\small,
  breaklines=true,
  frame=single,
  numbers=left,
  numberstyle=\tiny,
  keywordstyle=\color{blue},
  commentstyle=\color{gray},
}

\pagestyle{fancy}
\fancyhf{}
\fancyhead[L]{IF4020 Kriptografi}
\fancyhead[R]{Tugas 2: \algname}
\fancyfoot[C]{\thepage}

% Foto anggota; jika file belum ada, tampil kotak placeholder
\newcommand{\foto}[1]{%
  \IfFileExists{figures/#1}{\includegraphics[width=3cm,height=4cm,keepaspectratio]{#1}}%
  {\fbox{\parbox[c][4cm][c]{3cm}{\centering\small Foto\\(#1)}}}}

% Sel anggota di cover: nama, NIM, file foto
\newcommand{\anggotacell}[3]{%
  \begin{minipage}[t]{0.45\textwidth}\centering
    \foto{#3}\\[2pt] {\bfseries #1}\\ #2
  \end{minipage}}

% Blok tanda tangan: nama, NIM
\newcommand{\ttd}[2]{%
  \begin{minipage}[t]{0.45\textwidth}\centering
    \vspace{2.5cm}\rule{0.8\linewidth}{0.4pt}\\ #1 (#2)
  \end{minipage}}
